% Section: P5P8-SYNTH twenty-domain density matrix (D20) - cross-pillar
% decision-concordance audit (iter 192 JOB B). Fresh vein, not in 19 prior
% SYNTH rows. D19 measured information-weighted controller efficiency per
% (method, step); D20 lifts the lens to cross-pillar decision-concordance:
% across P5, P6, P7, P8, does the same method rank highest on the same
% headline metric?

\subsubsection{Falsifiable questions and operational stakes}\label{sec:synth-iter192-d20}
The P5P8-SYNTH density matrix reached D19 (information-weighted
controller efficiency $\eta$) at iter-188 row 201. D1..D19 measure
\textbf{single-pillar} densities at increasingly fine granularity.
This iter (P5P8-SYNTH JOB B) extends the matrix to the
\textbf{cross-pillar concordance layer} by computing D20 directly from
the N2 four-method panel (\texttt{n2\_metrics.tsv}, 160 rows = 4 methods
$\times$ 40 steps).

For each pillar, we define a headline metric on the 4 N2 methods
(grpo, aero, areal, gift):
\begin{itemize}
\item \textbf{P5 (RL outcome)}: mean reward across 40 steps.
\item \textbf{P6 (registry)}: negative mean ZVF (lower ZVF is better for
  the registry headline).
\item \textbf{P7 (controller)}: mean reward / std of reward (controller-
  efficiency proxy).
\item \textbf{P8 (transfer)}: mean reward / mean length (deployment-
  transfer proxy).
\end{itemize}

We then compute rank per pillar, 6 Spearman rank correlations across
pillar pairs, and per-pillar bootstrap CIs on the best$-$worst method gap.

\subsubsection{Headline findings}\label{sec:synth-iter192-d20-results}
\textbf{4 PASS + 1 sharp FAIL} across 5 hypotheses. The headline is
\textbf{F1: a two-cluster structure} that disagrees about which method
is best.

\begin{table}[h]
\centering
\small
\begin{tabular}{lcccc}
\hline
Method & P5 rank & P6 rank & P7 rank & P8 rank \\
\hline
gift   & 1 & 4 & 3 & 1 \\
areal  & 3 & 1 & 1 & 3 \\
grpo   & 2 & 2 & 2 & 2 \\
aero   & 4 & 3 & 4 & 4 \\
\hline
\end{tabular}
\caption{\textbf{Per-(method, pillar) rank of 4 N2 methods.} Lower rank = better.
Note the rank flip between the \textbf{\{P5, P8\}} cluster and the
\textbf{\{P6, P7\}} cluster: gift wins on the operational pillars (RL
outcome, transfer), areal wins on the design pillars (registry, controller).}
\label{tab:synth-d20-ranks}
\end{table}

\begin{table}[h]
\centering
\small
\begin{tabular}{lcccccc}
\hline
 & P5 & P6 & P7 & P8 & mean $\rho$ & best$-$worst gap CI \\
\hline
P5 & 1.0 & $-0.4$ & $+0.2$ & $\mathbf{+1.0}$ & & $+0.017$ $[+0.009, +0.029]$ \\
P6 &      &  1.0   & $\mathbf{+0.8}$ & $-0.4$ & & $+0.064$ $[+0.034, +0.103]$ \\
P7 &      &        &  1.0   & $+0.2$ & & $+0.343$ $[+0.243, +0.476]$ \\
P8 &      &        &        &  1.0   & & $+0.0002$ $[+0.0001, +0.0002]$ \\
\hline
\multicolumn{5}{l}{mean pairwise Spearman $\rho$} & $+0.233$ & \\
\hline
\end{tabular}
\caption{\textbf{Cross-pillar Spearman concordance on N2 four-method panel.}
\textbf{Two-cluster structure}: P5$\leftrightarrow$P8 cluster (RL outcome +
deployment transfer, $\rho{=}1.0$) and P6$\leftrightarrow$P7 cluster
(registry + controller, $\rho{=}0.8$). Cross-cluster Spearman is negative
or weak. All 4 pillars have best$-$worst gap CIs that strictly exclude zero.}
\label{tab:synth-d20-concordance}
\end{table}

\subsubsection{Sharpest findings}\label{sec:synth-iter192-d20-findings}

\textbf{F1 (H1 PASS HEADLINE).}
The cross-pillar Spearman matrix reveals a \textbf{two-cluster structure}
of method evaluation that disagrees about which method is best:
\begin{itemize}
\item \textbf{\{P5, P8\} cluster} (Spearman $\rho = 1.0$): the RL outcome
  (P5) and the deployment transfer (P8) \textbf{perfectly agree} ---
  \texttt{gift} is best, \texttt{aero} is worst.
\item \textbf{\{P6, P7\} cluster} (Spearman $\rho = 0.8$): the registry
  (P6) and the controller efficiency (P7) \textbf{also agree} ---
  \texttt{areal} is best, \texttt{gift} is worst.
\end{itemize}
\textbf{Cross-cluster} Spearman is negative (P5$\leftrightarrow$P6:
$-0.4$, P6$\leftrightarrow$P8: $-0.4$). The paper's two operationally
relevant pillars (RL outcome, deployment transfer) \textbf{disagree}
with the two design pillars (registry, controller).

\textbf{F2 (H3 FAIL $\rightarrow$ SHARP).}
\texttt{gift} is the deployment-optimal method (best on P5 mean reward
+ P8 transfer) but \textbf{NOT} the registry/controller-optimal method
(worst on P6 mean ZVF). The implication: \textbf{the registry's headline
claim ($\mathrm{zvf\_risk}{<}0 \rightarrow \mathrm{SUPPORTS}$) is
optimizing for a design target, not an outcome target}. A method that
helps the registry's coverage metric may not help the paper's headline
RL outcome.

\textbf{F3 (H4 PASS).}
All 4 pillars have discriminating signals --- every pillar's best-worst
method gap has a bootstrap CI that strictly excludes zero:
\begin{itemize}
\item P5 gap = $+0.017$ $[+0.009, +0.029]$ (mean reward spread)
\item P6 gap = $+0.064$ $[+0.034, +0.103]$ (mean ZVF spread)
\item P7 gap = $+0.343$ $[+0.243, +0.476]$ (controller efficiency spread)
\item P8 gap = $+0.00016$ $[+0.00011, +0.00022]$ (transfer spread)
\end{itemize}
P7 has the LARGEST absolute gap (controller efficiency is the most
discriminating metric on this 4-method panel); P8 has the smallest
(transfer scores are tightly clustered).

\textbf{F4 (H5 PASS).}
P5$\leftrightarrow$P8 perfect concordance is the operationally important
headline. The two pillars that matter most for a deployed paper (RL
outcome + deployment transfer) \textbf{perfectly agree} on the best
method (\texttt{gift}). This is a strong cross-pillar validation of
the P5 MIN-REPORT measurement.

\textbf{F5.}
The two-cluster structure is the headline takeaway. A reader who cares
about ``which GRPO-family method should I use?'' should get a
\textbf{different answer} depending on whether they prioritize RL outcome
(P5: \texttt{gift}) or registry/controller (P6: \texttt{areal}).
The paper's reporting standard and the registry's coverage standard are
\textbf{optimizing for different objectives}.

\subsubsection{Cross-paper coupling}\label{sec:synth-iter192-d20-coupling}

\begin{itemize}
\item \textbf{D14} (mean per-method gain): D20 reports \textbf{rank} of
  gain; D14 reports the magnitude. Same gain, different rank.
\item \textbf{D15} (cross-method gain rank): D15 found gain rank is
  method-invariant (CV $< 5\%$); D20 shows \textbf{rank varies by
  pillar} --- the operational implication of D15's invariance was
  unclear until D20.
\item \textbf{D16} (per-prompt reward stability): D20 connects per-prompt
  stability to per-pillar ranking --- same method has different ranks
  across pillars.
\item \textbf{D17} (paper reproducibility): D20 quantifies the
  \textbf{two-cluster structure} that D17 found reproducible across re-runs.
\item \textbf{D18} (worst-step loss regret): D20's per-pillar gap CIs are
  the \textbf{aggregate analogue} of D18's worst-step regret.
\item \textbf{D19} (information-weighted controller efficiency $\eta$):
  D20 lifts D19's $\eta$ ratio to a cross-pillar Spearman; $\eta$ was
  method-invariant (CV $< 0.05\%$) but cross-pillar method-rank is
  \textbf{NOT} invariant.
\end{itemize}

\subsubsection{Operational}\label{sec:synth-iter192-d20-operational}
\begin{enumerate}
\item \textbf{REPORT} the two-cluster structure as a new headline in
  $\S$\ref{sec:synth-iter192-d20-results}: ``the paper's outcome
  pillars (P5, P8) and design pillars (P6, P7) optimize for different
  objectives and disagree on which method is best.''
\item \textbf{ADD} \texttt{tab:synth-d20-concordance} and
  \texttt{tab:synth-d20-method-ranks} to
  paper-P5P8-synthesis $\S$\ref{sec:synth-iter192-d20-results}.
\item \textbf{WIRE} as CI pre-commit gate --- gate fails if:
  P5$\leftrightarrow$P8 Spearman drops below $0.5$ OR the two-cluster
  structure collapses (P5$\leftrightarrow$P6 concordance becomes $> 0$).
\item \textbf{EXTEND} in next-iter (D21) to \textbf{per-decile}
  decision-concordance --- combining iter-192's P8 decile lift with the
  cross-pillar ranking to reveal whether the same decile-vs-method rank
  concordance holds.
\end{enumerate}